\uuid{6sLr}
\titre{Distance}
\auteur{bodin ; hanani ; mzari}
\organisation{exo7}
\licence{CC BY-NC-SA 4.0}
\datecreate{2020-01-20}
\niveau{L1}
\module{Géométrie}
\chapitre{Géométrie du plan}
\sousChapitre{Distance, norme, produit scalaire}
\difficulte{1}
\theme{distance}
\exo7qcmid{lille-1.geometrie-plane.001}

\contenu{
\begin{question}
\qclassification{140.01,140.02}
\qtags{motcle=distance}
On considère les points $A(3,0)$ et $B(0,4)$. Quelle est la distance $d$ entre $A$ et $B$ ?
\begin{answers}
    \bad{$d=3$}
    \bad{$d=4$}
    \good{$d=5$}
    \bad{$d=3+4=7$}
\end{answers}
\begin{explanations}
D'abord, $\overrightarrow{AB}=(-3,4)$. Donc $d=\sqrt{(-3)^2+4^2}=\sqrt{25}=5$.
\end{explanations}
\end{question}
}
